% 7. Zaključak
\section{Zaključak}
\label{sec:zakljucak}

U ovom radu predstavljen je EduGrader, višemodelska platforma otvorenog koda za ocenjivanje pitanja otvorenog tipa uz podršku veštačke inteligencije i pod nadzorom nastavnika.
Sistem objedinjuje više velikih jezičkih modela (GPT-4o, GPT-4o-mini, DeepSeek-Chat i dva modela usmerena na rezonovanje, GPT-5.1 i DeepSeek-Reasoner) u podesiv tok ocenjivanja koji podržava evaluaciju i u PRA i u GRA režimu.
EduGrader je evaluiran na 686 autentičnih studentskih odgovora sa šest univerzitetskih kurseva na dve institucije, koji obuhvataju kratke činjeničke odgovore, teorijska objašnjenja, izraze nalik kodu i dodatni skup dužih odgovora esejskog tipa.

Empirijski rezultati pokazuju da PRA režim dosledno daje višu usklađenost sa ljudskim ocenjivačima od GRA režima, kroz sve modele i konfiguracije strogosti.
Modeli usmereni na rezonovanje, naročito GPT-5.1, postižu najjače slaganje sa ljudskim ocenjivanjem, dostižući Pirsonov koeficijent korelacije od 0{,}90 u neutralnom PRA režimu.
Ovi ishodi povoljno se porede sa prethodno objavljenim rezultatima u literaturi i sugerišu da, za evaluirane tipove odgovora --- a posebno za relativno strukturirana STEM pitanja --- savremeni veliki jezički modeli mogu aproksimirati obrasce ljudskog ocenjivanja sa odstupanjima uporedivim sa prirodnom ljudskom varijabilnošću.
Istovremeno, analize otkrivaju sistematske razlike između modela, režima ocenjivanja i konfiguracija strogosti, što naglašava važnost pažljivog izbora modela i parametara pri stvarnoj primeni.

Iako EduGrader donosi značajne dobitke u efikasnosti, podesivu strogost i detaljno generisanje povratnih informacija, rezultati potvrđuju da su sistemi veštačke inteligencije najefikasniji kada se koriste kao pomoćni alati, a ne kao zamena za čoveka.
Ljudski nadzor ostaje neophodan za pravednost, kontekstualno tumačenje i rešavanje dvosmislenih ili graničnih slučajeva.

Objavljivanjem sistema EduGrader kao projekta otvorenog koda nastoji se podstaći transparentnost, reproducibilnost i zajedničko unapređivanje sistema.
